% The PACS viewer in full: what it shows (cines at their frame rate, stacks with window and level,
% structured reports, PDF, videos, pictures), the files it reads, the controls and keys, the
% calipers, opening files and folders, the series list, the themes, the addresses, dicom-py
% (serve, index, wfdb), the examples.
% Sources: viewers/pacs/README.md, viewers/pacs/python/dicom/README.md,
% viewers/pacs/examples/README.md, checked against viewers/pacs/index.html, app.js, dicom.js,
% style.css and viewers/pacs/python/dicom/dicom_index_py/cli.py and page_server.py.
\chapter{The PACS viewer}
\label{chap:pacs-viewer}

This chapter describes the PACS viewer: what it shows, the files it reads, its controls, and
how to run it alone or with \cmd{dicom-py}. It is written for clinicians and researchers who
review images, and for administrators who deploy the viewer.

The PACS viewer shows cardiology and other DICOM studies (echo and other ultrasound loops,
angiography (XA), CT, MR, nuclear medicine and X-ray), structured reports and PDF documents, and
pictures and videos. It is one of Medorama's viewers, and also a page of its own that needs no
installation. Open it on the website (\url{https://medorama.org/viewers/pacs/}),
or open \file{viewers/pacs/index.html} from a clone, from disk or from any web server. Then choose
\ui{Open record\dots} (a study's or series' files, a video, a picture, a PDF) or
\ui{Open multiple records\dots} (a folder: every study in it and in its subfolders, listed by
series), or drop files or a folder on the page. The files are read and drawn in the browser, and
sent nowhere. The page has a light and a dark theme (\cref{sec:pacs-viewer-themes}).

\begin{note}[For research]
The PACS viewer is for research, not for diagnosis: a PACS used to make diagnoses is a regulated
medical device.
\end{note}

\screenshot{echo.png}{The PACS viewer in the dark theme, as Medorama shows it (the bar above it is Medorama's, with a synthetic patient): a synthetic echocardiography loop, apical four-chamber view, playing at 30 frames per second; the series list on the left.}{fig:pacs-viewer-echo}

% ==========================================================================================
\section{What it shows}
\label{sec:pacs-viewer-shows}

The files opened together become a list of items, each shown in its own way. When the files come
from several studies, the list groups them by study.

\begin{xltabular}{\linewidth}{@{}P{2.0cm}P{5.0cm}L@{}}
\caption{The items of the PACS viewer.}
\label{tab:pacs-viewer-items}\\
\toprule
{\tabhead Item} & {\tabhead From} & {\tabhead Shown} \\
\midrule
\endfirsthead
\toprule
{\tabhead Item} & {\tabhead From} & {\tabhead Shown} \\
\midrule
\endhead
\bottomrule
\endfoot
Cine loop & A multi-frame DICOM file of ultrasound, angiography, fluoroscopy, IVUS, OCT or
  endoscopy, or any multi-frame file with a frame time or a cine rate; a video (MP4, M4V, WebM,
  MOV, OGV); pictures given a frame rate & Played at its own rate (\code{FrameTime},
  \code{CineRate}, the frame time vector; \code{fps=} for pictures), in a loop, frame by frame,
  at $\frac{1}{4}\times$ to $2\times$ \\
Stack & The single-frame DICOM images of one series (CT, MR, NM, X-ray), or a multi-frame file
  that is not a cine & Ordered along the slice direction (else by instance number); the wheel or
  \key{$\uparrow$} \key{$\downarrow$} move through it, and it can play too \\
Still & One image or picture (PNG, JPEG, GIF, WebP, BMP) & As it is \\
Report & A DICOM structured report (echo, cath, others) & Its measurements with their units, and
  its text, in its sections; never the observer's or the patient's name, nor the subject
  context (the subject's name, ID, UID, birth date; the accession number) \\
Document & A DICOM encapsulated PDF, a PDF file & In the browser's PDF viewer; a DICOM document
  whose MIME type is not \code{application/pdf} is shown elsewhere \\
Elsewhere & A DICOM ECG or Holter waveform; a DICOM document that is not a PDF; a DICOM file
  whose compression the browser does not decode & Said, in place of the image
  (\cref{sec:pacs-viewer-elsewhere}) \\
\end{xltabular}

\screenshot{pacs.png}{A synthetic CT series (24 slices) as a stack, in the PACS viewer inside Medorama (dark theme): the series in the list, the slice and its position at the top right, the window and the level at the bottom left.}{fig:pacs-viewer-stack}

\subsection{Reading the images}
\label{sec:pacs-viewer-reading}

\begin{itemize}
\item \textbf{Window and level.} A drag on a gray image sets the window: to the right, a wider
  window; down, a higher level. The presets of \ui{Window} set it at once; a double click resets
  it (and the zoom and the position). Without a preset, the window is the file's, else one
  computed from the first frame, kept for the whole item so that a cine does not flicker.
\item Rescale slope and intercept are applied (CT in Hounsfield units), MONOCHROME1 is inverted,
  non-square pixels are drawn to scale, and 8-bit images (display-ready: ultrasound,
  angiography) show their whole range.
\item \textbf{Zoom and pan.} \key{Ctrl} + wheel (\key{Cmd} + wheel on macOS) zooms, and so do
  \ui{$-$} and \ui{+}, from 25\,\% to 800\,\%; \key{Shift} + drag pans, and so does a drag with
  the middle button, or on a color image, a picture or a video.
\item \textbf{A cine} plays as soon as it is chosen. \key{Space} plays or pauses it,
  \key{$\leftarrow$} \key{$\rightarrow$} step one frame, \key{$\uparrow$} \key{$\downarrow$} go to the
  previous or the next item.
\item \textbf{A stack} opens at its middle image; the wheel, \key{$\uparrow$} \key{$\downarrow$} (or
  \key{$\leftarrow$} \key{$\rightarrow$}, \key{PgUp} \key{PgDn}) move through it, \key{Home} and
  \key{End} go to its first and last image, and \key{Space} plays it.
\end{itemize}

\begin{table}[htbp]
\centering
\caption{The window presets (for gray images).}
\label{tab:pacs-viewer-presets}
\begin{tabularx}{\linewidth}{@{}P{3.0cm}P{2.6cm}L@{}}
\toprule
\tabhead Preset & \tabhead Width & \tabhead Level \\
\midrule
\ui{Auto} & The file's, or computed & The file's, or computed \\
\ui{Soft tissue} & 400 & 40 \\
\ui{Lung} & 1500 & $-600$ \\
\ui{Bone} & 1800 & 400 \\
\ui{Brain} & 80 & 40 \\
\ui{Liver} & 160 & 60 \\
\ui{Cardiac} & 600 & 300 \\
\bottomrule
\end{tabularx}
\end{table}

After a drag, the list shows \ui{Custom}. \ui{Window} is grayed out for color images, pictures
and videos.

\subsection{The corners}
\label{sec:pacs-viewer-corners}

The corners of the image show the series, the frame or slice, the window and the size, never the
patient:
\begin{itemize}
\item top left: the series' description, the modality, the slice thickness and the heart rate,
  when the file gives them;
\item top right: \code{Frame 17 / 40} (a cine) or \code{Image 13 / 24} (a stack), and a slice's
  location in millimeters;
\item bottom left: the window and the level of a gray image (\code{W 400 $\cdot$ L 40});
\item bottom right: the image's size in pixels, the zoom when it is not 100\,\%, and a cine's
  frame rate times the speed (\code{160 $\times$ 120 $\cdot$ 30 fps}).
\end{itemize}

\subsection{Elsewhere: what the page does not show}
\label{sec:pacs-viewer-elsewhere}

\begin{itemize}
\item A DICOM ECG waveform is said to open in the ECG viewer (\code{A DICOM ECG waveform: it opens
  in the ECG viewer.}), which reads it (\cref{chap:ecg-viewer}).
\item A DICOM Holter recording is said to be read another way: VitalStream does not read DICOM, so
  the recording is converted to a WFDB record first, with \cmd{dicom-py wfdb}
  (\cref{sec:pacs-viewer-dicom-py}), and that record opened in VitalStream (\code{A DICOM Holter
  recording: VitalStream does not read DICOM. Convert it to a WFDB record (dicom-py wfdb FILE.dcm
  OUT/) and open that record in VitalStream.}).
\item A DICOM file in JPEG-LS, JPEG 2000 or a video transfer syntax (MPEG-2, MPEG-4, HEVC) names
  its compression and says that browsers do not decode it, then how to read the file another way:
  \code{To read the file, open it in a desktop DICOM viewer, or convert it to an uncompressed
  transfer syntax first (for example gdcmconv --raw IN.dcm OUT.dcm; a video's frames with
  ffmpeg).}
\item A DICOM encapsulated document is shown only when its MIME type is \code{application/pdf},
  and always as a PDF; a document of another type (HTML, XML, a 3D model) is said instead
  (\code{A document of type TYPE: the viewer shows PDF documents only. Open it in a desktop
  DICOM viewer.}).
\end{itemize}

% ==========================================================================================
\section{The files it reads}
\label{sec:pacs-viewer-files-read}

DICOM Part 10 files, or bare datasets, in explicit or implicit VR little endian. Pixels may be
native (8, 16 or 32 bits, signed or not, gray or RGB, planar or not), RLE, or JPEG:

\begin{table}[htbp]
\centering
\caption{The JPEG transfer syntaxes the PACS viewer decodes.}
\label{tab:pacs-viewer-jpeg}
\begin{tabularx}{\linewidth}{@{}LP{4.6cm}@{}}
\toprule
\tabhead JPEG & \tabhead Decoded by \\
\midrule
Baseline and extended, 8 bits & The browser \\
Extended, 12 bits (CT) & \file{jpeg.js} \\
Lossless (process 14, any predictor, 2 to 16 bits) & \file{jpeg.js} \\
\bottomrule
\end{tabularx}
\end{table}

A file has no more frames than its pixel data hold: a \code{NumberOfFrames} larger than that (a
damaged or crafted header) is cut to the frames present.

JPEG-LS, JPEG 2000 and the video transfer syntaxes (MPEG-2, MPEG-4, HEVC) are named but not
decoded in the browser: the viewer says so, and the file is read another way, in a desktop DICOM
viewer or converted first to an uncompressed transfer syntax, such as with \cmd{gdcmconv -{}-raw}
(\cref{sec:pacs-viewer-elsewhere}). Big-endian and deflated files are not read
(\code{big-endian DICOM files are not read here}, \code{deflated DICOM files are not read here}).
A folder's \file{DICOMDIR}, its text files (\file{.txt}, \file{.md}, \file{.csv}, \file{.json},
\file{.xml}, \file{.html}) and its hidden files are passed over. Besides DICOM, the viewer reads
videos (MP4, M4V, WebM, MOV, OGV, as far as the browser plays them), pictures (PNG, JPEG, GIF,
WebP, BMP) and PDF files; any other file is taken for DICOM.

A file that cannot be read is listed with its reason when nothing else can be shown
(\code{These files are not ones the viewer reads.}, with the first five reasons, such as
\code{NAME: not a DICOM file} or \code{NAME: a DICOM object without an image}); among other files,
it is passed over. A frame that cannot be decoded (a damaged file) is said in place of the image
(\code{This frame cannot be decoded: \dots}), and the cine stops; the other items stay usable.

% ==========================================================================================
\section{The page}
\label{sec:pacs-viewer-page}

The header holds, from left to right: the name \ui{PACS} and the item's title (or the
\code{name=} of the address), the controls of the image, and the buttons to open files. Below it,
the series list on the left, and the image, with its corners and a slider under it.

\begin{xltabular}{\linewidth}{@{}P{4.6cm}L@{}}
\caption{The controls of the PACS viewer.}
\label{tab:pacs-viewer-controls}\\
\toprule
{\tabhead Control} & {\tabhead Does} \\
\midrule
\endfirsthead
\toprule
{\tabhead Control} & {\tabhead Does} \\
\midrule
\endhead
\bottomrule
\endfoot
\ui{\textless} (\key{$\leftarrow$}) & The previous frame or image \\
The play button (\key{Space}) & Play or pause a cine, a stack or a video \\
\ui{\textgreater} (\key{$\rightarrow$}) & The next frame or image \\
The speed: \ui{$\frac{1}{4}\times$}, \ui{$\frac{1}{2}\times$}, \ui{1$\times$}, \ui{2$\times$} & The
  playing speed, times the item's frame rate \\
\ui{Window} & The window presets of gray images (\cref{tab:pacs-viewer-presets}) \\
\ui{$-$} \ui{+} (\key{Ctrl} + wheel) & Smaller, larger \\
\ui{Calipers} & Measure a distance (\cref{sec:pacs-viewer-calipers}) \\
\ui{Server folders\dots} & The folders of the computer serving the page (\cmd{dicom-py serve}
  only) \\
\ui{Open record\dots} & Files: a study's or series' DICOM files, a video, a picture or a PDF \\
\ui{Open multiple records\dots} & A folder: every study in it and in its subfolders \\
The moon or sun button & Switches between the light and the dark theme \\
The series list & The items, with a thumbnail, a title and a line (the modality, the series
  number, and the number of images, or of frames and the frame rate, or \code{report},
  \code{document}, \code{video}); under \ui{Study 1}, \ui{Study 2}, \dots{} when there are several
  studies. A click shows an item \\
The slider & The images or frames of the item \\
\end{xltabular}

The controls that do not apply to an item are grayed out: the frame controls for a single image,
the window for color images, the zoom and the calipers for a report or a document.

\begin{xltabular}{\linewidth}{@{}P{5.2cm}L@{}}
\caption{The keys and mouse actions of the PACS viewer.}
\label{tab:pacs-viewer-keys}\\
\toprule
{\tabhead Action} & {\tabhead Effect} \\
\midrule
\endfirsthead
\toprule
{\tabhead Action} & {\tabhead Effect} \\
\midrule
\endhead
\bottomrule
\endfoot
\key{Space} & Play or pause \\
\key{Esc} & Clear the calipers' points \\
\key{$\leftarrow$} \key{$\rightarrow$} on a cine or a video & The previous, the next frame (a
  video: a thirtieth of a second, paused) \\
\key{$\uparrow$} \key{$\downarrow$} on a cine or a video & The previous, the next item of the list \\
\key{$\uparrow$} \key{$\leftarrow$} \key{PgUp}, \key{$\downarrow$} \key{$\rightarrow$} \key{PgDn} on a
  stack & The previous, the next image \\
\key{Home}, \key{End} on a stack & The first, the last image \\
Wheel over the image & The previous or the next frame or image \\
\key{Ctrl} (or \key{Cmd}) + wheel & Zoom \\
Drag on a gray image & The window: right wider, down a higher level \\
\key{Shift} + drag; a drag with the middle button; a drag on a color image, picture or video &
  Pan \\
Double click & Reset the window (\ui{Auto}), the zoom and the position \\
With \ui{Calipers} on, a click & A point of the calipers \\
\end{xltabular}

The keys work when the focus is not in a list or a text box.

\subsection{Calipers}
\label{sec:pacs-viewer-calipers}

\ui{Calipers} turns the calipers on (the button is highlighted) or off. Click two points (a click
also stops a playing cine): the viewer draws a dashed line between them and writes the distance.
A third click starts again; \key{Esc} clears the points.

\begin{table}[htbp]
\centering
\caption{The unit of the calipers' distance.}
\label{tab:pacs-viewer-calipers}
\begin{tabularx}{\linewidth}{@{}P{2.0cm}L@{}}
\toprule
\tabhead Distance in & \tabhead When \\
\midrule
Cm & The ultrasound region's scale (\code{PhysicalDeltaX}, \code{PhysicalDeltaY}) at the first
  point \\
Mm & The pixel spacing (\code{PixelSpacing}); for projection images, \code{ImagerPixelSpacing}
  (marked \code{at the detector}) \\
Mm & Pictures and videos opened with \code{mm=} (millimeters per pixel) \\
Pixels & Otherwise \\
\bottomrule
\end{tabularx}
\end{table}

\subsection{Light and dark}
\label{sec:pacs-viewer-themes}

The PACS viewer has a light and a dark theme, and is dark until told otherwise, as reading
rooms usually are. The moon button (in the light theme) or the sun button (in the dark theme)
switches the theme, and the choice is remembered in the browser.
\code{theme=light} or \code{theme=dark} in the address sets it for that page, without remembering
it. The image itself is always drawn on black.

% ==========================================================================================
\section{Opening studies}
\label{sec:pacs-viewer-open}

\begin{table}[htbp]
\centering
\caption{The ways to open files in the PACS viewer.}
\label{tab:pacs-viewer-open}
\begin{tabularx}{\linewidth}{@{}P{4.2cm}L@{}}
\toprule
& \tabhead Opens \\
\midrule
\ui{Open record\dots} & The files chosen: a study's or a series' DICOM files, a video, a picture
  or a PDF \\
\ui{Open multiple records\dots} & A folder: every file in it and in its subfolders, grouped by
  study and listed by series \\
Drop & The files, or the folder, dropped on the page (a folder: every file in it and in its
  subfolders, hidden ones passed over) \\
\ui{Server folders\dots} & Shown only when the page is served by \cmd{dicom-py serve}: the folders
  of the serving computer (\cref{sec:pacs-viewer-serve}) \\
\bottomrule
\end{tabularx}
\end{table}

Files and folders whose names start with \code{.} are passed over, among the files chosen and in
a folder dropped: a dropped folder is walked through all its subfolders, every hidden file and
folder left out, and each folder's entries are read in full (the browser gives them in batches,
read until none is left). Every file is held in memory at
once, so at most 5000 files and 4\,GB are read together; more is refused, with the advice to
open one study's folder (\code{This folder holds N files (X GB): more than 5000 files or 4 GB,
which are held in memory at once. Open one study's folder.}).

\paragraph{How the list is made.} The files are read in the order of their names. Single-frame
images of the same series form one stack (a stack of one image is a still); each multi-frame
file is a cine or a stack of its own; each report, document, video and picture is an item. Pictures
opened with \code{fps=} in the address, two or more, are one cine instead, in the order of their
names. The items are listed by study, in the order the studies were met, then by series number,
then by instance number, then by title. The first item is shown.

% ==========================================================================================
\section{Running and deploying}
\label{sec:pacs-viewer-run}

\begin{xltabular}{\linewidth}{@{}P{5.0cm}L@{}}
\caption{The ways to run the PACS viewer.}
\label{tab:pacs-viewer-run}\\
\toprule
{\tabhead How} & {\tabhead What it opens} \\
\midrule
\endfirsthead
\toprule
{\tabhead How} & {\tabhead What it opens} \\
\midrule
\endhead
\bottomrule
\endfoot
The website, \url{https://medorama.org/viewers/pacs/} & The files you
  pick or drop, read in your browser: they never leave your computer \\
\file{viewers/pacs/index.html} from a clone, opened from disk (\code{file://}) & The same,
  offline \\
Any web server: copy the page's files (\file{index.html}, \file{style.css} and the \file{.js}
  files of \file{viewers/pacs}) & The same; and files on that server, by address (\code{\#src=},
  \code{\#list=}) \\
\cmd{dicom-py serve} on the computer where the studies are & The same; and the folders of that
  computer, under a root you choose (\ui{Server folders\dots}) \\
Inside Medorama (\cmd{tools/install-medorama.sh}) & A patient's studies, as a tab of the
  patient's page, de-identified by Medorama's server (\cref{chap:viewers-in-medorama}) \\
\end{xltabular}

The page needs no build and no dependencies; it uses nothing but the browser. For a quick look
from a web server: \cmd{python3 -m http.server -{}-directory viewers/pacs}.

\subsection{Served from a computer: \texorpdfstring{\cmd{dicom-py serve}}{dicom-py serve}}
\label{sec:pacs-viewer-serve}

The Python package \file{viewers/pacs/python/dicom} (\code{dicom\_index\_py}, the command
\cmd{dicom-py}) serves the page together with the folders of the computer it runs on, so that
studies on a server are read without copying them, through an SSH tunnel to its port (the
commands to install it on each operating system are at \url{https://medorama.org/docs/install.html}). \cmd{serve}
uses only the Python standard library; its module, \file{page\_server.py}, is the same file as
\cmd{muse-py serve}'s.

\begin{lstlisting}[style=shell]
pip install ./viewers/pacs/python/dicom        # copies the page into the package
dicom-py serve                                 # http://127.0.0.1:8081/browse: your home folder
dicom-py serve --root /data/echo               # another folder
dicom-py serve --set-password                  # once, on a shared machine (VitalStream's password)
\end{lstlisting}

\begin{lstlisting}[style=signature]
dicom-py serve [--root DIR] [--host HOST] [--port N] [--auth-file FILE]
               [--set-password | --clear-password] [--no-auth] [--no-browser]
\end{lstlisting}

\begin{params}
\item[--root DIR] The folder served: only the files and folders under it are listed and served,
  every symbolic link resolved (a link that leads outside is refused, and not listed). Default:
  the home folder. It lies inside the root the package was installed with: a folder that
  contains that root is narrowed to it, one outside it is refused.
\item[--host HOST] Where it listens. Default: \code{127.0.0.1}. A request must name the server
  by a loopback name, an IP address or this host; any other name is refused (no DNS rebinding).
\item[--port N] Default: 8081.
\item[--auth-file FILE] The password file, shared with VitalStream and \cmd{muse-py serve}.
  Default: \file{\textasciitilde/.config/vitalstream/auth.json}.
\item[--set-password] Set (or change) the password the page asks for, then exit (at least 8
  characters, asked twice). For \cmd{muse-py serve} and \cmd{dicom-py serve}, a sign-in lasts 30
  days, until the server stops, or until a browser signs out, which ends every sign-in; after a
  wrong password the next try, from any browser, waits 1\,s, then 2, 4 \dots{} up to 60\,s.
\item[--clear-password] Remove the password, then exit.
\item[--no-auth] Run without asking for the password.
\item[--no-browser] Do not open a browser window.
\end{params}

\paragraph{The server folders.} \file{browse} (titled \ui{PACS viewer: server folders}) lists the
folders of a folder under the root, and the files the viewer reads: DICOM files, with the
extension \file{.dcm} or \file{.dicom} or without an extension, videos (\file{.mp4},
\file{.webm}, \file{.mov}), pictures (\file{.png}, \file{.jpg}, \file{.jpeg}) and PDFs, with their
size. Hidden entries are neither listed nor served, and a folder shows its first 5000 entries. A
click on a file opens it in the viewer (\code{\#src=files/PATH}); \ui{Open the records in this
folder} opens every file of the folder (\code{\#list=list/PATH/}, at most 5000). The browser's
Back button returns to the list. A file or a folder opened from \file{browse} is shown with
\code{\#src=} (or \code{\#list=}), where the page hides its \ui{Server folders\dots} button; the
page opened at its own address (\url{http://127.0.0.1:8081/}) shows \ui{Server folders\dots},
which leads to the list, when \code{GET served} answers 200 \code{\{"viewer", "root"\}}
(the requests of \cmd{serve}: \cref{sec:ecg-viewer-serve}). A folder's listing says when it
holds more than 5000 entries (\code{"truncated": true}).

Files are sent as they are read, in part when a browser asks for a range (\code{Range}), under a
policy that runs nothing in them (\code{Content-Security-Policy: sandbox}); a kind of file the
viewers do not read is sent as a download. The sign-in cookie is signed with a key of the
running server too: browsers send a cookie to every port of a host, and another server on
127.0.0.1 cannot use it.

\paragraph{Fixed at install time.} As for \cmd{muse-py}, two variables of the environment are
fixed in the package when it is installed: \code{VIEWER\_ROOT}, the most \cmd{serve} ever serves
(\code{\textasciitilde} allowed; default the whole file system: a \opt{-{}-root} outside it is
refused, one that contains it is narrowed to it), and \code{VIEWER\_PASSWORD=0}, for no password
support. An editable install (\cmd{pip install -e}) serves the page from the clone; the variables,
when given, are written into the clone (a file git ignores), and without them its root is the
whole file system, with password support.

\begin{lstlisting}[style=shell]
VIEWER_ROOT=/data/echo VIEWER_PASSWORD=1 pip install ./viewers/pacs/python/dicom
\end{lstlisting}

\paragraph{On a server where the studies are:} install with a \code{VIEWER\_ROOT} that holds them,
set a password once (\cmd{dicom-py serve -{}-set-password}), and start it without a browser, in
\cmd{tmux} or as a service: \cmd{dicom-py serve -{}-root /data/echo -{}-no-browser}. From your
computer: \cmd{ssh -N -L 8081:127.0.0.1:8081 user@server}, then open
\url{http://127.0.0.1:8081/}. Without a password, anyone who can connect to the port can read the
files under the root (\cref{chap:viewers-alone}).

% ==========================================================================================
\section{dicom-py: databases of DICOM studies}
\label{sec:pacs-viewer-dicom-py}

\cmd{dicom-py} also builds a database to import into Medorama from a folder of DICOM files, and
converts DICOM ECGs to WFDB records. It uses pydicom and NumPy, which \cmd{pip} installs with it.

\begin{lstlisting}[style=shell]
dicom-py index /path/to/dicom -o imaging.csv   # a database: a row per loop, series, report
dicom-py wfdb ECG.dcm out/                     # a DICOM ECG (12-lead or Holter) as a WFDB record
medorama import-table imaging.csv              # then: the rows reach their patients
dicom-py --version
\end{lstlisting}

\apientry{dicom-py index}{A database of DICOM studies}
\begin{lstlisting}[style=signature]
dicom-py index FOLDER -o CSV
\end{lstlisting}
Reads every DICOM file under \file{FOLDER} (hidden files and folders passed over), headers only,
never the pixels, and writes a CSV with one row per object a viewer shows; a file it cannot read
(a damaged header, sequences nested too deep) is listed on the standard error, and the others
are indexed. It prints \code{CSV: N object(s)}, and \code{; M not linked}
with each of them, and the reason, on the standard error.
\begin{params}
\item[-o, --output CSV] The CSV to write (required); the file paths in it are relative to its
  folder.
\end{params}

\begin{xltabular}{\linewidth}{@{}P{4.4cm}L@{}}
\caption{The rows of \cmd{dicom-py index}.}
\label{tab:pacs-viewer-index-rows}\\
\toprule
{\tabhead Object} & {\tabhead Row} \\
\midrule
\endfirsthead
\toprule
{\tabhead Object} & {\tabhead Row} \\
\midrule
\endhead
\bottomrule
\endfoot
An ultrasound loop, an angiography run, any multi-frame file & One per file, with its frames and
  frame rate \\
A series of single-frame images (CT, MR, NM) & One per series, linking the series' folder: only
  when that folder holds one patient's files (else it is listed as not linked, and no row is
  written) \\
A single image (an X-ray) & One per file \\
A structured report (echo, cath \dots) & One per file; each measurement a column
  \code{Name (unit)}, such as \code{Left Ventricular Ejection Fraction (\%)} (a repeated
  measurement keeps its first value, or its mean when one is given); its text in
  \code{interpretation} \\
An encapsulated PDF & One per file \\
A DICOM ECG (12-lead, Holter) & One per file; its measurements as the ECG viewer names them
  (\code{QT (ms)}, \code{QRS (ms)}, \code{Vent. rate (bpm)} \dots), its statements in
  \code{interpretation} \\
\end{xltabular}

The columns are \code{mrn} (the patient ID\phimark), \code{time} (the device's clock),
\code{category} and \code{name}, then \code{modality}, \code{description}, \code{study\_uid},
\code{series\_uid}, \code{images}, \code{frames} and \code{fps}; then the measurement columns,
then \code{interpretation} and \code{file} (relative to the CSV). The categories are \code{echo}
(ultrasound), \code{cath} (angiography, fluoroscopy, IVUS, OCT), \code{ecg} and \code{holter} (ECG
waveforms), and \code{img} (everything else); a report or a PDF takes its category from its
title's words (echo, cath, ECG). Columns of different databases are never merged by Medorama: an
echo database's \code{LVEF (\%)} and a cath database's are distinct (\cref{chap:ecg-imaging}).

\begin{phibox}
The patient ID is written, because it is the key that links each row to its patient. Keep the
database as PHI, with Medorama's databases, never in the repository. Names and dates of birth are
never written.
\end{phibox}

\apientry{dicom-py wfdb}{A DICOM ECG as a WFDB record}
\begin{lstlisting}[style=signature]
dicom-py wfdb FILE OUT_DIR [--name NAME]
\end{lstlisting}
Converts the DICOM ECG \file{FILE} (12-lead or Holter) into a WFDB record in \file{OUT\_DIR}, and
prints its path: format 16, at 1\,$\mu$V per unit, with the annotations' measurements and
statements in the header's comments (each value on one line; each line of a text annotation a
statement). The record opens in the ECG viewer and, for a Holter
recording, in VitalStream.
\begin{params}
\item[--name NAME] The record's name. Default: the file's name without its extension.
\end{params}

% ==========================================================================================
\section{Addresses}
\label{sec:pacs-viewer-address}

Another program opens the PACS viewer at a file, or a folder, by its address's fragment. The
files are fetched with the page's own credentials; changing the fragment shows other files
without reloading the page.

\begin{table}[htbp]
\centering
\caption{The fragments of the PACS viewer's address.}
\label{tab:pacs-viewer-address}
\begin{tabularx}{\linewidth}{@{}P{4.6cm}L@{}}
\toprule
\tabhead Fragment & \tabhead Means \\
\midrule
\code{src=URL} & One file \\
\code{list=URL} & A folder: \code{URL} answers \code{\{"files": [\dots]\}}, and each file is read
  beside it \\
\code{mm=0.2} & The millimeters per pixel of pictures and videos (the calipers) \\
\code{fps=30} & The frame rate of pictures, which are then played as one loop \\
\code{name=\dots} & The name to show \\
\code{theme=light} or \code{dark} & The theme \\
\code{who=\dots}, \code{what=\dots}, \code{pos=\dots}, \code{nav=\dots} & A host's context, shown in
  place of the file's name (\cref{chap:viewers-alone}) \\
\bottomrule
\end{tabularx}
\end{table}

Opened at a file or a folder, the page hides its buttons that open files, \ui{Server
folders\dots} and the theme button: the host decides what is shown. Medorama fills \code{mm} and
\code{fps} from the columns of the row that linked the file (\code{pixel spacing (mm)},
\code{frame rate (fps)}, \dots).

% ==========================================================================================
\section{The examples}
\label{sec:pacs-viewer-examples}

\file{viewers/pacs/examples/} holds synthetic DICOM objects to try the viewer with: every value
and every image is invented, and no file describes a real person. They are written by
\file{tools/make-viewer-examples.py} with the demo's generators. On the website, the examples page
(\url{https://medorama.org/viewers/pacs/examples/}) has a link that opens
each file in the live viewer; from a clone, open \file{viewers/pacs/index.html} and choose a file
with \ui{Open record\dots}, or the whole folder with \ui{Open multiple records\dots}.

\begin{table}[htbp]
\centering
\caption{The PACS viewer's examples.}
\label{tab:pacs-viewer-examples}
\begin{tabularx}{\linewidth}{@{}P{3.2cm}L@{}}
\toprule
\tabhead File & \tabhead What \\
\midrule
\file{echo\_a4c.dcm} & An echocardiography loop (ultrasound, apical four-chamber view, 30
  frames): a dilated, poorly contracting left ventricle \\
\file{angio\_lca.dcm} & A coronary angiography run (XA) of the left coronary artery \\
\file{chest\_xray.dcm} & A frontal chest radiograph (DX) with a right lower-lobe opacity \\
\bottomrule
\end{tabularx}
\end{table}

% ==========================================================================================
\section{Files and tests}
\label{sec:pacs-viewer-files}

\begin{table}[htbp]
\centering
\caption{The files of the PACS viewer (\file{viewers/pacs}).}
\label{tab:pacs-viewer-files}
\begin{tabularx}{\linewidth}{@{}P{3.6cm}L@{}}
\toprule
\tabhead File & \tabhead Contents \\
\midrule
\file{dicom.js} & The DICOM reader: elements, sequences, pixel data (native, RLE, JPEG: the
  browser's or \file{jpeg.js}), windows, frame rate, ultrasound regions, structured reports,
  encapsulated documents, waveforms; \file{viewers/ecg/dicom.js} is the same file (the tests
  check it) \\
\file{jpeg.js} & Lossless JPEG and 12-bit DCT JPEG, which browsers do not decode (no
  dependencies) \\
\file{host.js} & A host's context (\code{who}, \code{what}, \code{pos}, \code{nav}), shown in place
  of the file's name; the same file in each viewer \\
\file{app.js}, \file{index.html}, \file{style.css} & The page \\
\file{python/dicom/} & \cmd{dicom-py}: databases of DICOM studies, DICOM ECGs as WFDB (pydicom),
  and \cmd{serve} \\
\file{examples/} & Synthetic DICOM objects to try the viewer with \\
\file{tests/} & The tests of the page and of \cmd{dicom-py} \\
\bottomrule
\end{tabularx}
\end{table}

\begin{lstlisting}[style=shell]
node --test viewers/pacs/tests/*.test.js     # the page's reader (DICOM files made in the tests)
pytest viewers/pacs/python/dicom/tests       # dicom-py, serve (its root and its password)
\end{lstlisting}
